\chapter{Natural Gas and LNG}\label{ch:m3:natural-gas-and-lng}

In early 2020 a megawatt-hour of gas for next-month delivery in the Netherlands cost
well under 25 euros. Two and a half years later it traded above 300. In the same
months a tanker could load gas in Louisiana at a price linked to the US benchmark,
a few dollars per million British thermal units, sail for two weeks, and sell the
cargo into Europe at many times that. Gas is a network commodity: it moves in
pipelines, and where pipelines end, in ships that must first cool it to
$-\qty{162}{\degreeCelsius}$. Its prices are local, seasonal and, when the network
breaks, extreme. This chapter covers the hubs where gas is priced, the basis between
them, the seasonal cycle of storage, the arbitrage of liquefied cargoes, and the crisis
of 2022.

\section{Gas as a network commodity: hubs and units}

Gas is traded at hubs, points of a pipeline network at which title can pass
between parties. Some hubs are physical places; in Europe most are notional.

\begin{definition}[Virtual trading point, Henry Hub, Title Transfer Facility]\label{def:m3:natural-gas-and-lng:hubs}
A \emph{virtual trading point}\index{virtual trading point} is a notional point of a
gas transmission system, operated by its network operator, at which gas that has
entered the system can be traded without regard to where it physically is.
\emph{Henry Hub}\index{Henry Hub} is a physical interconnection of pipelines in
Louisiana whose price is the US benchmark and the delivery point of the main US gas
futures. The \emph{Title Transfer Facility}\index{Title Transfer Facility} (TTF) is the
virtual trading point of the Dutch network and the most traded European gas benchmark.
\end{definition}

Units differ by region. The United States prices gas in dollars per million British
thermal units (MMBtu); continental Europe in euros per megawatt-hour of energy; the
United Kingdom in pence per therm (\qty{0.1}{\mmbtu}). One MMBtu is \qty{0.293071}{MWh}:
a TTF price of \qty{40}{EUR/MWh} at EURUSD 1.16 is \qty{13.60}{\$/\mmbtu}. The US
futures contract is 10{,}000 MMBtu delivered at \omterm{def:m3:natural-gas-and-lng:hubs}{Henry Hub}; its trading ends three
business days before the first day of the delivery month.

\begin{definition}[Gas day]\label{def:m3:natural-gas-and-lng:gasday}
The \emph{gas day}\index{gas day} is the 24-hour period over which a network balances
flows and settles daily products. In the European Union it runs from 05:00 to 05:00
UTC in winter time and from 04:00 to 04:00 UTC in summer time.
\end{definition}

\section{Locational basis and pipelines}

A price at a hub is a price for gas there. Between two hubs the price difference
is bounded by the cost of moving gas between them when capacity is free, and
unbounded when it is not.

\begin{definition}[Locational basis]\label{def:m3:natural-gas-and-lng:basis}
The \emph{locational basis}\index{locational basis} of a hub is its price minus the
price of a reference hub (in the United States, \omterm{def:m3:natural-gas-and-lng:hubs}{Henry Hub}) for the same delivery
period. It is traded as basis swaps and futures.
\end{definition}

\begin{proposition}[Basis is bounded by transport only when capacity is available]\label{prop:m3:natural-gas-and-lng:basis}
Let $c$ be the variable cost of moving gas from hub $A$ to hub $B$ and suppose a
shipper holds unused capacity. Then $P_B - P_A \le c$: otherwise the shipper buys at
$A$, ships and sells at $B$ for a riskless profit. Without unused capacity the
difference can be any size, and its value accrues to the holders of capacity.
\end{proposition}

\begin{proof}
With capacity, the trade $-P_A - c + P_B > 0$ is available to anyone holding it and
is done until the prices move. Without capacity no trade links them.
\end{proof}

A producing region whose pipelines out are full can see its hub price fall far below
the reference, sometimes below zero; a consuming region cut off from supply can see
it rise far above. The value of pipeline capacity is the expected value of these
spreads, which is why capacity is itself sold at auction.

\section{Storage and seasonality}

Demand for gas peaks in winter (heating) and, in hot regions, in summer (power for
air conditioning), while production is roughly flat. Storage bridges the two.

\begin{definition}[Injection season, summer--winter spread]\label{def:m3:natural-gas-and-lng:season}
The \emph{injection season}\index{injection season} is the part of the year in which
storage is filled, roughly April to October in the northern hemisphere; gas is
withdrawn in the winter months that follow. The \emph{summer--winter
spread}\index{summer--winter spread} is the price of gas for the next winter minus
the price for the preceding summer; it pays for storing gas through the season.
\end{definition}

\begin{omfigure}
\centering
\begin{tikzpicture}[font=\small]
\begin{axis}[width=11cm,height=4.6cm,
  xlabel={delivery month},ylabel={forward price (stylised)},
  tick label style={font=\small},label style={font=\small},
  xmin=0.5,xmax=12.5,ymin=0,ymax=20,grid=major,grid style={gray!25},
  xtick={1,2,3,4,5,6,7,8,9,10,11,12},
  xticklabels={Apr,May,Jun,Jul,Aug,Sep,Oct,Nov,Dec,Jan,Feb,Mar},
  yticklabels={}]
\addplot[omDef,very thick,mark=*,mark size=1.4pt] coordinates {(1,8.6) (2,8.3) (3,8.2) (4,8.3) (5,8.4) (6,8.6) (7,9.0) (8,10.2) (9,11.2) (10,11.6) (11,11.3) (12,10.4)};
\draw[omThm,thick,{Stealth[length=2mm]}-{Stealth[length=2mm]}] (axis cs:3,8.2) -- (axis cs:3,11.6);
\node[omThm,font=\footnotesize,anchor=north west,align=left,fill=white,inner sep=1pt] at (axis cs:0.7,19.2) {summer--winter spread: pays for\\storing gas from summer to winter};
\draw[gray,dashed] (axis cs:3,11.6) -- (axis cs:10,11.6);
\node[font=\footnotesize,anchor=south] at (axis cs:4,1.2) {injection};
\node[font=\footnotesize,anchor=south] at (axis cs:10,1.2) {withdrawal};
\end{axis}
\end{tikzpicture}

\omcaption{The seasonal shape of a gas forward curve, schematic: summer months are
cheap, winter months dear. A storage owner buys summer and sells winter; the spread,
net of injection and withdrawal costs, is its intrinsic value. Stylised, no
data.}\label{fig:m3:natural-gas-and-lng:season}
\end{omfigure}

The simplest valuation of storage buys the cheapest summer months and sells the
dearest winter ones forward, within the facility's capacity and its injection and
withdrawal rates; the rest of its value lies in re-optimising as prices move, which
One Quant Book 6, chapter 16, values as options.

\begin{dated}{2026-09}{European storage rules}\label{dat:m3:natural-gas-and-lng:storage}
Regulation (EU) 2022/1032 required member states' storage to be 80\% full by 1 November
2022 and 90\% by 1 November of each later year. Regulation (EU) 2025/1733, published on
10 September 2025, extends the rules to the end of 2027 and lets the 90\% target be met
at any time between 1 October and 1 December, with deviations of up to 10 percentage
points in difficult market conditions.
\end{dated}

\section{LNG: liquefaction, shipping and cargo arbitrage}

\begin{definition}[Liquefied natural gas]\label{def:m3:natural-gas-and-lng:lng}
\emph{Liquefied natural gas}\index{liquefied natural gas} (LNG) is natural gas cooled
to about $-\qty{162}{\degreeCelsius}$, at which it is a liquid about six hundred times
denser than the gas: it is shipped in insulated carriers, loses a small share to
boil-off on the way, and is warmed back to gas (regasified) at the receiving terminal.
\end{definition}

LNG turns regional gas markets into a global one, but only at the margin, because the
chain is expensive (\cref{fig:m3:natural-gas-and-lng:chain}). Its contracts come in two
families. Long-term contracts to Asian buyers were long priced on oil.

\begin{definition}[Oil indexation, destination flexibility]\label{def:m3:natural-gas-and-lng:contracts}
\emph{Oil indexation}\index{oil indexation} prices gas by a formula $P = s \cdot P_{\mathrm{oil}}
+ k$, a slope $s$ times an average of past oil prices plus a constant. A contract has
\emph{destination flexibility}\index{destination flexibility} when the buyer may take the
cargo to any destination rather than to a named terminal.
\end{definition}

The US export model is different: a buyer pays a fixed fee per MMBtu for the right to
liquefaction capacity, whether or not it lifts cargoes, plus a variable price of 115\% of
\omterm{def:m3:natural-gas-and-lng:hubs}{Henry Hub} for each cargo it lifts, \omterm{def:m3:physical-commodity-markets:incoterms}{free on board} at the terminal; it may cancel a cargo
with notice and still owes the fixed fee. Such a buyer holds a portfolio of options: each month it lifts the cargo
if the best \omterm{def:m3:physical-commodity-markets:netback}{netback} covers the variable price, and sends it to the destination that
pays the most.

\begin{definition}[Japan Korea Marker]\label{def:m3:natural-gas-and-lng:jkm}
The \emph{Japan Korea Marker}\index{Japan Korea Marker} (JKM) is a \omterm{def:m3:physical-commodity-markets:pra}{price-reporting
agency}'s daily assessment of spot LNG cargoes delivered ex-ship into Japan, South Korea,
China and Taiwan, the main Asian spot benchmark.
\end{definition}

\begin{omfigure}
\centering
\begin{tikzpicture}[font=\small,
  box/.style={draw,thick,rounded corners=2pt,align=center,minimum height=1cm,text width=2.2cm},
  arr/.style={-{Stealth[length=2.2mm]},thick}]
\node[box,draw=omDef,fill=omDef!8] (hh) at (0,0) {pipeline gas\\at Henry Hub};
\node[box,draw=omProp,fill=omProp!8] (liq) at (3.2,0) {liquefaction\\(FOB)};
\node[box,draw=omExo,fill=omExo!8] (ship) at (6.4,0) {carrier\\voyage};
\node[box,draw=omProp,fill=omProp!8] (reg) at (9.6,0) {regasification\\terminal};
\draw[arr] (hh) -- (liq);
\draw[arr] (liq) -- (ship);
\draw[arr] (ship) -- (reg);
\node[font=\footnotesize,align=center,anchor=north] at (3.2,-0.75) {$1.15\times$ Henry Hub\\+ fixed fee};
\node[font=\footnotesize,align=center,anchor=north] at (6.4,-0.75) {charter, fuel,\\boil-off};
\node[font=\footnotesize,align=center,anchor=north] at (9.6,-0.75) {regas fee;\\sold at the hub};
\end{tikzpicture}

\omcaption{The LNG chain of a US export cargo, with the cost each link adds. The \omterm{def:m3:physical-commodity-markets:netback}{netback}
is the destination price less the right-hand costs; the cargo is lifted if the \omterm{def:m3:physical-commodity-markets:netback}{netback}
covers the variable price at liquefaction. Schematic.}\label{fig:m3:natural-gas-and-lng:chain}
\end{omfigure}

\begin{example}[Europe or Asia]\label{ex:m3:natural-gas-and-lng:choice}
Take a cargo of 3.5 million MMBtu, a charter rate of \$60{,}000 a day, boil-off of
0.1\% a day, and one-way voyages of 14 days to north-west Europe and 25 to north Asia
(illustrative, but close to the 14 days from Corpus Christi to Rotterdam and 28 to Hong Kong
through Panama that the Oxford Institute for Energy Studies gives), with regasification at \qty{0.50}{} and \qty{0.40}{\$/\mmbtu}. With
\omterm{def:m3:natural-gas-and-lng:hubs}{Henry Hub} at \qty{3.00}{} the variable price is \qty{3.45}{\$/\mmbtu}; shipping costs
\qty{0.54}{} to Europe and \qty{0.97}{\$/\mmbtu} to Asia, so Asia must pay
\qty{0.33}{\$/\mmbtu} more than Europe to win the cargo. At \qty{12.00}{} in Europe and
\qty{12.50}{} in Asia the \omterm{def:m3:physical-commodity-markets:netback}{netbacks} are \qty{10.96}{} and \qty{11.13}{\$/\mmbtu}: Asia.
\end{example}

\begin{omfigure}
\begin{tikzpicture}
\begin{axis}[width=11cm,height=5cm,ymode=log,
  ylabel={\$/MMBtu (log scale)},
  tick label style={font=\small},label style={font=\small},
  xmin=2000,xmax=2026.75,ymin=1,ymax=100,grid=major,grid style={gray!25},
  ytick={2,5,10,20,50,100},yticklabels={2,5,10,20,50,100},
  xtick={2000,2004,2008,2012,2016,2020,2024},xticklabel style={/pgf/number format/1000 sep={}},
  legend columns=3,legend style={font=\footnotesize,at={(0.5,-0.18)},anchor=north,draw=none,fill=none,/tikz/every even column/.append style={column sep=8pt}},
  table/col sep=comma]
\addplot[omDef,thick] table[x=t,y=hh]{figdata/markets-3/04-natural-gas-and-lng/hubs.csv};
\addplot[omThm,thick] table[x=t,y=eu]{figdata/markets-3/04-natural-gas-and-lng/hubs.csv};
\addplot[omProp,thick,dashed] table[x=t,y=jp]{figdata/markets-3/04-natural-gas-and-lng/hubs.csv};
\legend{Henry Hub,European hub,Asian LNG}
\end{axis}
\end{tikzpicture}

\omcaption{Monthly averages of three gas benchmarks, January 2000 to July 2026, on a log
scale. US gas decoupled from the rest of the world after 2009; Europe peaked at
\qty{69.98}{\$/\mmbtu} in August 2022. Data: FRED series MHHNGSP (US Energy Information
Administration), PNGASEUUSDM and PNGASJPUSDM (International Monetary
Fund).}\label{fig:m3:natural-gas-and-lng:hubs}
\end{omfigure}

Run month by month since the first US Gulf export cargoes of 2016, the \omterm{def:m3:physical-commodity-markets:netback}{netback} rule of
\cref{ex:m3:natural-gas-and-lng:choice} sends the cargo to Asia in 93 of 125 months and
says it should not be lifted in seven months of 2020 (February to August) and in
February 2021, when a freeze in Texas sent \omterm{def:m3:natural-gas-and-lng:hubs}{Henry Hub} up
(\cref{fig:m3:natural-gas-and-lng:netback}). Buyers did cancel cargoes in 2020: the US
Energy Information Administration counted 45 cancelled for August and about 30 for
September, as low spot prices in Europe and Asia made US exports uneconomic.

\begin{omfigure}
\begin{tikzpicture}
\begin{axis}[width=11cm,height=4.8cm,ymode=log,
  ylabel={\$/MMBtu (log scale)},
  tick label style={font=\small},label style={font=\small},
  xmin=2016,xmax=2026.75,ymin=0.5,ymax=100,grid=major,grid style={gray!25},
  ytick={1,2,5,10,20,50,100},yticklabels={1,2,5,10,20,50,100},
  xtick={2016,2018,2020,2022,2024,2026},xticklabel style={/pgf/number format/1000 sep={}},
  legend columns=2,legend style={font=\footnotesize,at={(0.5,-0.18)},anchor=north,draw=none,fill=none,/tikz/every even column/.append style={column sep=8pt}},
  table/col sep=comma]
\addplot[omDef,very thick] table[x=t,y=netback]{figdata/markets-3/04-natural-gas-and-lng/netback.csv};
\addplot[omThm,thick,dashed] table[x=t,y=fobvar]{figdata/markets-3/04-natural-gas-and-lng/netback.csv};
\legend{best netback at the US Gulf,variable price $1.15\times$ Henry Hub}
\end{axis}
\end{tikzpicture}

\omcaption{A US Gulf cargo's best \omterm{def:m3:physical-commodity-markets:netback}{netback} against its variable price, monthly, March 2016
to July 2026. Where the \omterm{def:m3:physical-commodity-markets:netback}{netback} falls below the dashed line the cargo is worth more
cancelled. Illustrative shipping and regasification costs of
\cref{ex:m3:natural-gas-and-lng:choice}; prices as
\cref{fig:m3:natural-gas-and-lng:hubs}.}\label{fig:m3:natural-gas-and-lng:netback}
\end{omfigure}

\begin{dated}{2026-09}{US exports}\label{dat:m3:natural-gas-and-lng:us}
US LNG exports rose 26\% to 15.1 billion cubic feet a day in 2025, 26\% of the world
total; the United States is the largest LNG exporter, ahead of Qatar and Australia.
\end{dated}

\section{The crisis of 2022}

Russia's pipeline exports to the European Union and the United Kingdom fell by almost
40\% in the first seven months of 2022 against 2021; Nord Stream 1 ran at 20\% of its
capacity in July and stopped entirely at the start of September (in November Sweden's Security Service
found that detonations on both Nord Stream pipelines had been gross sabotage); in October Russian
pipeline deliveries were 80\% below their level a year earlier. European hub prices rose above 300 euros a megawatt-hour, the front month peaking in August.
Europe replaced the lost pipeline gas with LNG, outbidding Asia for flexible cargoes,
and filled its storage under the new rules. Utilities that had sold their production
forward on exchanges faced margin calls that ran to billions: the importer Uniper had drawn a EUR~9 billion state
credit line in full by the end of August (\cref{ch:m3:when-markets-break-together}).

\begin{dated}{2026-09}{The market correction mechanism}\label{dat:m3:natural-gas-and-lng:mcm}
Council Regulation (EU) 2022/2578 set a dynamic bidding limit on front-month TTF
derivatives: if the front-month settlement exceeded EUR~180/MWh and was EUR~35 above a
global LNG reference price, orders more than EUR~35/MWh above that reference were not
to be accepted. It applied from 1 February 2023 and, extended once, until 31 January 2025.
\end{dated}

\section{Tutorial: netbacks and the destination choice}\label{tut:m3:natural-gas-and-lng:netback}

\begin{tutorial}
\textbf{Goal.} Compute a cargo's \omterm{def:m3:physical-commodity-markets:netback}{netbacks}, choose its destination and decide whether to
lift it, month by month. \textbf{End state:}
\cref{ex:m3:natural-gas-and-lng:choice}, \cref{fig:m3:natural-gas-and-lng:netback} and
the count of months not lifted.
\begin{tutsteps}
\item \textbf{Shipping, \omterm{def:m3:physical-commodity-markets:netback}{netback}, choice.}
\omcode{code/firm/lngarb/firm_lngarb.py}{38}{57}{Shipping cost per delivered MMBtu, netback and the best destination.}\label{lst:m3:natural-gas-and-lng:netback}
\item \textbf{Break-even and lift.}
\omcode{code/firm/lngarb/firm_lngarb.py}{60}{68}{The break-even spread between two destinations and the lift rule.}\label{lst:m3:natural-gas-and-lng:lift}
\item \textbf{Run} \texttt{m3\_lng.history()}, \texttt{summary()} and \texttt{fig\_lng.py}.
\end{tutsteps}
\textbf{What to change next.} Route Asian cargoes around the Cape of Good Hope when the
Panama Canal is restricted, and count how many destination decisions change; make the
charter rate move with the spread, as it does when many cargoes chase the same arbitrage.
\end{tutorial}

\section{Build: the cargo arbitrage engine}\label{bld:m3:natural-gas-and-lng:lngarb}

\begin{build}
\bfield{Purpose} The miniature firm holds liquefaction capacity: every month it must
decide whether to lift and where to send each cargo, and it quotes gas across hubs in
three units.
\bfield{Interface} \texttt{eur\_mwh\_to\_usd\_mmbtu}, \texttt{usd\_mmbtu\_to\_eur\_mwh};
\texttt{fob\_price(henry\_hub, slope, fee)}; \texttt{Route(name, days, charter\_per\_day,
boiloff\_per\_day, regas)}; \texttt{shipping\_cost}; \texttt{netback}; \texttt{choose};
\texttt{breakeven\_spread}; \texttt{lift}.
\bfield{Rules} The fixed fee is sunk and never enters the lift decision; boil-off is paid
at the cargo's FOB value and reduces the delivered volume; the charter is paid for the
round trip; ties go to the alphabetically first destination.
\bfield{Acceptance tests} \texttt{code/firm/lngarb/tests/}: unit round trips; a longer
route costs more; the choice flips at the break-even spread; the lift rule.
\bfield{Stretch} Canal constraints and their waiting times; a fleet of carriers
scheduled across a year; the value of the lift option computed by Monte Carlo on
correlated hub prices.
\end{build}

\begin{omsources}
\item Regulation (EU) 2022/1032 and Regulation (EU) 2025/1733 on gas storage; Council
Regulation (EU) 2022/2578 (market correction mechanism); Commission Regulation (EU)
2015/703 (gas day).
\item US Energy Information Administration, \emph{Today in Energy}: LNG exports 2020 and
2025; Russia's pipeline exports to Europe (2022). IEA, \emph{2022 Energy Crisis: Frequently
Asked Questions}.
\item Cheniere Energy Partners, Form 10-K (SPA pricing at 115\% of Henry Hub plus a fixed
fee).
\item Swedish Security Service, ``Confirmed sabotage of the Nord Stream gas pipelines'', 18 November 2022 (Internet Archive
copy); Fortum, press release on Uniper's KfW facility, 29 August 2022; Oxford Institute for Energy Studies, \emph{LNG
Shipping Chokepoints}, NG 188, February 2024; ACER, preliminary report on the market correction
mechanism, 2023.
\item CME Group, Henry Hub natural gas futures overview; S\&P Global Platts, JKM.
\item D.~Dumitrescu et al., ``European energy crisis, the Dutch TTF, and the market
correction mechanism'', \emph{Journal of World Energy Law \& Business}, 2024.
\item FRED series MHHNGSP, PNGASEUUSDM, PNGASJPUSDM.
\end{omsources}

\section{Exercises}

\begin{exercise}[$\star$]\label{exo:m3:natural-gas-and-lng:1}
Convert \qty{10}{\$/\mmbtu} into EUR/MWh at EURUSD 1.16, and into pence per therm at
GBPUSD 1.30.
\end{exercise}

\begin{exercise}[$\star$]\label{exo:m3:natural-gas-and-lng:2}
Why is the \omterm{def:m3:natural-gas-and-lng:season}{summer--winter spread} positive in most years, and what bounds it?
\end{exercise}

\begin{exercise}[$\star$]\label{exo:m3:natural-gas-and-lng:3}
A US LNG buyer pays a fixed fee of \qty{2.50}{\$/\mmbtu} and 115\% of \omterm{def:m3:natural-gas-and-lng:hubs}{Henry Hub}. \omterm{def:m3:natural-gas-and-lng:hubs}{Henry Hub}
is \qty{2.00}{} and its best \omterm{def:m3:physical-commodity-markets:netback}{netback} \qty{2.10}{\$/\mmbtu}. Does it lift? What does it lose
per MMBtu of capacity either way?
\end{exercise}

\begin{exercise}[$\star\star$]\label{exo:m3:natural-gas-and-lng:4}
With the routes of \cref{ex:m3:natural-gas-and-lng:choice} and \omterm{def:m3:natural-gas-and-lng:hubs}{Henry Hub} at \$3, Europe
pays \qty{11.00}{\$/\mmbtu}. Above what Asian price does the cargo go to Asia?
\end{exercise}

\begin{exercise}[$\star\star$]\label{exo:m3:natural-gas-and-lng:5}
A hub's pipelines out are full and it trades \qty{2}{\$/\mmbtu} below \omterm{def:m3:natural-gas-and-lng:hubs}{Henry Hub} while
transport costs \qty{0.30}{}. Who earns the difference, and what trade would you do if you
held capacity?
\end{exercise}

\begin{exercise}[$\star\star$]\label{exo:m3:natural-gas-and-lng:6}
The voyage to Asia doubles to 50 days (a canal is closed). By how much does the break-even
spread rise, with the other inputs of \cref{ex:m3:natural-gas-and-lng:choice}?
\end{exercise}

\begin{exercise}[$\star\star\star$]\label{exo:m3:natural-gas-and-lng:7}
\emph{Coding.} With \texttt{history}, count the months since March 2016 in which the
cargo goes to Asia and list those in which it is not lifted.
\end{exercise}

\begin{exercise}[$\star\star\star$]\label{exo:m3:natural-gas-and-lng:8}
\emph{Find the flaw.} ``European gas was ten times dearer than US gas in August 2022, so
every US cargo earned ten times its cost.''
\end{exercise}

\section{Problem: Where Should the Cargo Go?}

\begin{problem}[{Weekend problem --- a flexible cargo between two continents}]\label{pb:m3:natural-gas-and-lng:1}
A trading house holds a US Gulf cargo of 3.5 million MMBtu with full \omterm{def:m3:natural-gas-and-lng:contracts}{destination
flexibility}, bought at 115\% of \omterm{def:m3:natural-gas-and-lng:hubs}{Henry Hub} plus a sunk fee. Use the routes and costs of
\cref{ex:m3:natural-gas-and-lng:choice}.

\medskip\noindent\textbf{Part I --- The costs.}
\begin{enumerate}
\item With \omterm{def:m3:natural-gas-and-lng:hubs}{Henry Hub} at \$3, give the variable price of the cargo.
\item Give the shipping cost per delivered MMBtu to Europe.
\item And to Asia.
\item What share of the cargo boils off on each voyage?
\item Why is the charter paid for twice the one-way days?
\end{enumerate}

\medskip\noindent\textbf{Part II --- The decision.}
\begin{enumerate}[resume]
\item Give the break-even Asia-minus-Europe spread.
\item With Europe at \$12.00 and Asia at \$12.50, where does the cargo go?
\item What is the cargo's margin over its variable price, in dollars?
\item What happens to the decision if the charter rate doubles?
\item In August 2022 (\omterm{def:m3:natural-gas-and-lng:hubs}{Henry Hub} \$8.81, Europe \$69.98, Asia \$54.16), where does it go?
\end{enumerate}

\medskip\noindent\textbf{Part III --- The option.}
\begin{enumerate}[resume]
\item Why is the right to choose a destination worth money even when prices are equal on
average?
\item In June 2020, what did the \omterm{def:m3:physical-commodity-markets:netback}{netback} rule say, and why?
\item What does the sunk fee change in that decision?
\item Who gains when the buyer cancels a cargo?
\item What market moves when many flexible cargoes chase the same arbitrage?
\end{enumerate}

\medskip\noindent\textbf{Part IV --- Judgement.}
\begin{enumerate}[resume]
\item Why did Europe win most flexible cargoes in 2022?
\item Why does a canal closure raise the price of gas in Asia relative to Europe?
\item Why might a buyer sign an oil-indexed contract rather than a Henry Hub-linked one?
\item State the \emph{named result}: the spread at which the cargo is indifferent between
the two destinations.
\item In one sentence: what links the world's gas prices?
\end{enumerate}
\end{problem}

\section{Interview questions}

\begin{interviewq}[$\star$ \iqroles{trader}]\label{iq:m3:natural-gas-and-lng:1}
What is the difference between \omterm{def:m3:natural-gas-and-lng:hubs}{Henry Hub} and TTF as hubs?
\end{interviewq}

\begin{interviewq}[$\star$ \iqroles{trader, researcher}]\label{iq:m3:natural-gas-and-lng:2}
Why is gas seasonal and oil much less so?
\end{interviewq}

\begin{interviewq}[$\star\star$ \iqroles{researcher}]\label{iq:m3:natural-gas-and-lng:3}
A US LNG offtake contract is a strip of options. On what, with what strike?
\end{interviewq}

\begin{interviewq}[$\star\star$ \iqroles{trader, risk}]\label{iq:m3:natural-gas-and-lng:4}
You are long gas at a hub and short \omterm{def:m3:natural-gas-and-lng:hubs}{Henry Hub}. What can make you lose a lot in a day?
\end{interviewq}

\begin{interviewq}[$\star\star$ \iqroles{risk}]\label{iq:m3:natural-gas-and-lng:5}
Why did European utilities that were fully hedged face a liquidity crisis in 2022?
\end{interviewq}

\begin{interviewq}[$\star\star\star$ \iqroles{developer, trader}]\label{iq:m3:natural-gas-and-lng:6}
Design the tool that tells an LNG desk every morning where each of its twenty cargoes
should go.
\end{interviewq}
